A Feed-Forward Neural Network (\textit{FFNN}) is our second custom implemented regression method. In this section we will explain the basics of how a FFNN works along with how we implemented these methods.

\subsection{The Structure}
A FFNN is a collection of nodes organized in layers. First, the input layer. The input layer number of nodes is equal to the number of features in the dataset. Second, the hidden layers. There can be an arbitrary number of hidden layers and of nodes per hidden layer. Generally, the more hidden layers or nodes per layer the more complex the Neural Network becomes, which makes it more prone to overfitting and more computationally heavy to train and run predictions on. Lastly, the output layer. The number of nodes in the output layer depends on the task, but for regression there is just one node, which give the value predicted value for the data points. 

\begin{figure} [h]
    \centering
    \caption{Structure of a FFNN}
    \label{fig:FFNN-structure}
 \end{figure}

\subsection{The Flow}
A FFNN at its core follows a very simple flow, see ~\autoref{fig:FFNN-flowchart}

 \begin{figure} [h]
    \centering
    \caption{Flowchart of a FFNN}
    \label{fig:FFNN-flowchart}
 \end{figure}

\subsection{Initialization}
We initialize an array of weights ($W_{ij}$) and an array of biases ($b_j$) for each pair of subsequent layers. The shape of $W_ij$ is the number of nodes in the previous layer ($i$) by the number of nodes in the current layer ($j$) which is also the shape of the bias array.\\
\\
To fill the arrays with initial values we use the He Initialization \cite{he2015delvingdeeprectifierssurpassing} by drawing from a zero-mean normal distribution with standard deviation of $\sqrt{2/n_i}$ which is considered to be among the best initializations for shallow ReLU based FFNNs.

\subsection{Forward Pass}
During the Forward Pass, the data, as the name suggests propagates forward through the network. At each layer, except the input layer, the activations of the previous layer ($a_i$) are weighed by the weights and have the bias added to them:
$$
z_j = a_i W_{ij} + b_j
$$
Then they are passed through the activation function to become the activation of the current layer, which are then passed to the next layer.
\subsubsection{Activation Functions}
There are many activation functions that one might choose depending on the task. Usually you use two different activation functions, one for the hidden layers and a different one for the output layer. For hidden layers we use the widely used activation function, the Rectified Linear Unit (\textit{ReLU}) function defined as $ReLU(x)=max(0,x)$. For the output layer of regression FFNNs a simple linear function is used $f(x) = x$
\subsection{Back propagation}

\subsection{Updating Parameters}

\subsection{Non-Essential Features}
\subsubsection{Early Stopping}
\subsubsection{Mini-Batch Gradient Descent}
\subsubsection{Weight Decay}